\chapter{Expenses and income}

\section{The problem with data entry}

Recording expenses is not a technical problem, it is a habit problem. Any
application that asks for six fields per coffee gets abandoned in a week.

The user of this application already had a habit: writing the month's expenses
in the notes app of a phone, as a list with bullets, sometimes with small sums
written out because three purchases went into the same category.

\begin{lstlisting}[style=sh]
June
- Expense: Groceries: 15-1+3
- Income: Salary: 1200
- Food: 8.53+1+26.25
- Transport: (10+5)*2
\end{lstlisting}

Rather than replace that habit, the application reads it. The user pastes the
text, and the parser turns it into rows.

\section{Parsing a pasted block}

The parser lives in \file{app/services/expenses\_parse.py} and works in three
steps.

\begin{enumerate}
  \item \textbf{Find the month.} If the first non-empty line is a month name and
        contains no colon, it is the title and it sets the month.
  \item \textbf{Clean each line.} A regular expression removes leading bullets:
        \code{-}, \code{*}, \code{\textbullet}, and a few others that phones
        insert.
  \item \textbf{Split by colon.} The last field is the arithmetic expression,
        everything before it is a label. If the first label is a kind
        (\emph{Expense} or \emph{Income}) it sets the kind, otherwise it is part
        of the category.
\end{enumerate}

Lines that do not fit come back in an \code{ignored} list instead of being
dropped in silence. The preview shows them, so a typo is visible rather than
mysterious.

\begin{lstlisting}[style=py]
@dataclass
class ParsedItem:
    kind: str          # expense | income
    category: str
    amount: Decimal
    note: str | None = None

@dataclass
class ParsedExpenses:
    title: str
    month: int | None
    items: list[ParsedItem]
    ignored: list[str] = field(default_factory=list)
\end{lstlisting}

\section{Evaluating arithmetic without eval}

A line can contain \code{15-1+3} or \code{(10+5)*2}. The obvious way to evaluate
that is \code{eval()}, and it is the wrong way: \code{eval} on text that came
from outside runs whatever that text says.

The parser uses Python's own \code{ast} module instead. The expression is parsed
into a syntax tree, and the tree is walked by a function that accepts only
numbers, unary signs, and six binary operators.

\begin{lstlisting}[style=py]
_BINOPS = {ast.Add: operator.add, ast.Sub: operator.sub,
           ast.Mult: operator.mul, ast.Div: operator.truediv,
           ast.FloorDiv: operator.floordiv, ast.Mod: operator.mod}

def _eval_node(node: ast.AST) -> float:
    if isinstance(node, ast.Constant):
        if isinstance(node.value, bool) or not isinstance(node.value, (int, float)):
            raise ExprError("numbers only")
        return float(node.value)
    if isinstance(node, ast.BinOp):
        op = _BINOPS.get(type(node.op))
        if op is None:
            raise ExprError("operator not allowed")     # ** lands here
        left, right = _eval_node(node.left), _eval_node(node.right)
        if op in (operator.truediv, operator.floordiv, operator.mod) and right == 0:
            raise ExprError("division by zero")
        return op(left, right)
    if isinstance(node, ast.UnaryOp) and isinstance(node.op, (ast.UAdd, ast.USub)):
        val = _eval_node(node.operand)
        return val if isinstance(node.op, ast.UAdd) else -val
    raise ExprError("expression not allowed")
\end{lstlisting}

Anything the walker does not recognise raises. A name, a function call, an
attribute or a subscript never reaches an evaluation branch, because there is no
branch for them.

\begin{pitfall}[Why exponentiation is blocked]
\code{**} is arithmetic and would fit the theme, but \code{2**9999999} is a
denial of service in eleven characters: it is a legal expression that takes a
long time and a lot of memory. It is left out of the operator table, and the
input is capped at 120 characters as a second line of defence.
\end{pitfall}

Two more guards run before parsing: the expression must match a character
allowlist (digits, spaces, \code{. , + - * / \% ( )}), and a comma is normalised
to a dot so that both decimal conventions work.

\section{The preview step}

Parsing produces a preview page, not database rows. The preview shows each
detected line with an editable category, an editable amount, a kind selector and
a checkbox, plus the list of ignored lines.

Saving offers two buttons, and the difference between them matters:

\begin{itemize}
  \item \textbf{Append} adds the rows to whatever the month already contains.
  \item \textbf{Replace} deletes every row of that month first.
\end{itemize}

Replace exists because the realistic workflow is pasting a month, noticing a
mistake, fixing the text and pasting again. Without it, the second paste would
double the month. It asks for confirmation in the browser, since it destroys
data.

\begin{lstlisting}[style=py]
if mode == "replace":
    db.execute(sql_delete(Transaction).where(
        Transaction.year == year, Transaction.month == month))
\end{lstlisting}

Note the two conditions: the delete is scoped to one month of one year, never to
the table. A test covers exactly that, by seeding a row in a different month and
asserting it survives.

\section{Reading amounts from a form}

Users type amounts in whatever convention they are used to. The form parser
accepts all four common forms:

\begin{lstlisting}[style=py]
def _parse_amount(raw: str) -> Decimal:
    s = raw.strip().replace(" ", "")
    if "," in s and "." in s:
        # the last separator is the decimal one; the other groups thousands
        if s.rfind(",") > s.rfind("."):
            s = s.replace(".", "").replace(",", ".")
        else:
            s = s.replace(",", "")
    else:
        s = s.replace(",", ".")
    try:
        return Decimal(s).quantize(Decimal("0.01"))
    except (InvalidOperation, ValueError):
        return Decimal("0.00")
\end{lstlisting}

The rule is simple and works for both conventions: whichever separator appears
last is the decimal separator. \code{1.234,56} and \code{1,234.56} both become
\code{1234.56}.

An unparseable value becomes zero rather than raising. A form that rejects the
whole submission because of one bad field is more annoying than a zero the user
can see and correct.

\section{The two views}

The monthly view lists the entries of one month with totals for income,
expenses and the balance, plus two doughnut charts by category. The yearly view
groups by month and shows a bar chart of income against expenses, plus expenses
by category for the whole year.

Both read through the same aggregation helpers used by the reports and the bot,
so the number on the page and the number in the PDF come from the same query.

\section{Investing is spending, and also not}

One category needed special treatment. Money moved into an investment leaves the
current account like an expense, but it is not consumption: it is still part of
net worth. Counting it as a plain expense makes a good month look terrible.

The application keeps it simple. Such entries are still expenses, but categories
named \emph{investment}, \emph{investments} or \emph{investing} are also summed
separately, and every summary adds a line: \emph{of which X was invested}.

\begin{lstlisting}[style=py]
INVESTMENT_CATEGORIES = {"investment", "investments", "investing"}

def invested_in_period(db, year, month) -> float:
    cond = [Transaction.year == year, Transaction.kind == "expense",
            func.lower(Transaction.category).in_(tuple(INVESTMENT_CATEGORIES))]
    if month:
        cond.append(Transaction.month == month)
    return _f(db.scalar(select(func.sum(Transaction.amount)).where(*cond)))
\end{lstlisting}

The comparison is done in lower case inside the query, so the user can capitalise
the category however they like.
